\section{结果与设计取舍}
\label{sec:results}

\subsection{两条路径在当前骨架上均有增量}
表\ref{tab:paths}报告两个预算截面的全部模型均值与种子标准差。
先行GEANT比较中，160轮预算下SPIN+Direct的NMAE为0.164810，
低于Context-only的0.168494与Direct-only的0.172246，分别改善2.19\%与4.32\%。
符合预定筛选规则后，原样扩展至Abilene，
对应数值为0.161358、0.164053和0.165011，分别改善1.64\%与2.21\%。
两个WAN在120轮截面也有相同排序。
更具体地，两个预算下，两个种子在三个条件中的NMAE均优于对应单路径。
这些结果补上了当前骨架的路径使用价值证据，
但不能把具有不同查询和解码设计的所有浅层SPIN实现一并排除。

固定降学习率阶段之后，三个模型的轨迹波动减小。
对每个WAN的两个种子及两个控制模型，共八个配对，
组合在第141--160轮的每个开发分数，均低于相应单路径在全160轮中的最佳分数。
因此，本轮排序并不只由某一个孤立低点产生。
这些是相关优化迭代的描述，既不是额外训练重复，也不是显著性证据；
统一日程同时增加了训练轮数，不能将全部变化单独归因于降低学习率。
完整逐轮曲线及逐种子、逐条件结果随复核材料保留。

\subsection{辅助指标与观测支持并非全面改善}
160轮预算下，组合的条件平均NRMSE在两个WAN和两个种子上，
均低于对应的两个单路径模型。
但GEANT相对Direct-only在不均匀和共享缺口条件的NRMSE，
两个种子均略差；均值改善主要得益于均匀条件。
因此，不能由条件平均值推导所有缺失模式均改善。

预定义支持分组进一步限定解释。
GEANT上，组合相对Context-only在两侧都有本流观测的内部位置改善，
但不均匀与共享缺口条件的单侧边界NMAE均值分别为
0.220937对0.220250、0.218896对0.217341，即存在小幅退步。
在受影响流的共同缺口内，三模型NMAE均值依次为
0.181171、0.183979和0.191145。
这些分组与局部读出提供插值偏置的解释相容，
不证明Direct恢复了SPIN必然丢失的幅值，也不构成新的主指标。

\subsection{增量成本与简单模型的取舍}
表\ref{tab:cost}显示，相对Context-only，加入Direct仅增加1,440个名义参数，
批量8每窗耗时在Abilene和GEANT分别增加3.82\%和3.47\%，
峰值已分配显存几乎不变。
这一结论不能扩大到全部执行情形：批量1延迟分别增加16.45\%和4.03\%。
同时，组合相对Direct-only的批量8耗时分别为11.79和16.45倍。
所以较低的NMAE并不意味着在吞吐受限系统中应一律承担上下文网络的成本。

历史Linear在相同开发流量与掩码上的NMAE分别为0.171021和0.167739，
本轮组合相对它分别改善5.65\%和1.75\%。
Abilene的NRMSE也改善9.35\%，但GEANT的NRMSE仍恶化0.42\%。
Linear的CPU成本已补测，并与GPU结果分别注明设备及计时边界，
不作跨设备加速比宣称。
旧四层SPIN和旧Direct+Sync参考保留在完整结果中，
其深度、训练日程或独立重复数量不同，不能据此完成公平外部方法优势证明。

\subsection{历史同步记忆对照的保留范围}
表\ref{tab:paths}同时保留原120轮Full结果。
在该历史配对中，增加同步记忆后，Abilene的NMAE基本持平且种子方向相反，
GEANT平均NMAE恶化1.41\%；条件平均NRMSE分别改善0.86\%和1.24\%。
历史同设备测量中，Full的批量8摊销成本分别增加约64\%和62\%。
这些事实解释了本轮优先验证无记忆组合的选择，
也保留了记忆对平方误差的收益。
本轮没有把Full延长至160轮，因此不据新旧预算交叉比较宣称同步机制普遍无效。
